\documentclass[11pt]{article}
\usepackage[margin=1in]{geometry}
\usepackage{amsmath,amssymb,amsthm,mathtools}
\usepackage{hyperref}
\setlength{\emergencystretch}{2em}
\newtheorem{theorem}{Theorem}
\newtheorem{lemma}[theorem]{Lemma}
\newtheorem{proposition}[theorem]{Proposition}
\newcommand{\C}{\mathbb C}
\newcommand{\Z}{\mathbb Z}
\newcommand{\R}{\mathbb R}
\newcommand{\ord}{\operatorname{ord}}
\newcommand{\lk}{\operatorname{lk}}
\title{From the actual M45 Seifert congruence to ambient equivalence\\
Unit-aware orders and the dimension-seven classification}
\author{Campaign 9000059: Engineer proof for independent review}
\date{November 2026}
\begin{document}
\maketitle

\begin{abstract}
We complete the geometric implication for the actual prepared germs of
the four/five manuscript, using the separately reviewed local-algebra,
integral tube, and integral-congruence results without alteration.
The Weierstrass units are retained in the computation of the two
initial forms. We check Kato's original smooth open-book theorem at
its inclusive dimension threshold, justify the actual compact pages'
handle condition and collar-fixed monodromy, and identify the ordered
filling form with the required linking form. The analytic ambient-pair
cone theorem then gives the homeomorphism of germs.
The accompanying exact computations certify the algebraic and
orientation calculations, not the geometric theorems. This is a prose
verification proof, not a formal-kernel certificate or a claim of
original discovery. Independent review of this final implication and
independent attribution of its literature inputs are still required.
The independent two/three manuscript is not addressed.
\end{abstract}

\section{Precise conclusion and retained inputs}

M45 is the staged September 27, 2026 manuscript
\emph{Ambiently homeomorphic isolated hypersurface germs in
\(\C^4\) with multiplicities four and five}.
Fix
\[
\begin{gathered}
r\geq4,\quad n_0>5\cdot2^r,\quad
n_0\equiv1\pmod{2^r},\quad n_0\equiv2\pmod5,\\
n_1=n_0+5\cdot2^r,\qquad
a=(n_0^2-1)/4,\qquad b=(n_1^2-1)/4.
\end{gathered}                                                    \tag{1}
\]
We retain the following separately reviewed results, including their
actual prepared germs, orientations, and parameter restrictions.
\begin{enumerate}
\item \texttt{C059BASE} and \texttt{C059APPLICATION}: the two finite
graph maps have degree four; the actual prepared functions have isolated
critical points, and
\(\mu=24ab+6a+6b+51\), whenever \(a,b>1\) and \(ab>5(a+b)\),
for every fixed \(s\ne0\).
\item \texttt{M45:thm:tube}: for the strict tube inequalities and a
sufficiently small common positive \(s\), the stated integral lattices
and their ordered filling forms are those of the actual local
Milnor pages, not merely rational models or groups of the same rank.
\item \texttt{M45:prop:arith-integral-congruence}: for (1) these actual
integral forms admit an onto integral isometry. This is an existence
theorem; no explicit enormous matrix is presumed.
\end{enumerate}
Their proofs and exact certificates remain in
\path{research/audit_report.tex},
\path{research/m45_seifert/}, and
\path{research/m45_arithmetic/}. They are not rerun or rewritten here.
The current ledger contains independent supporting judgements for
these inputs. Supporting judgements are not kernel verification.

\begin{theorem}[Verification of \texttt{M45:thm:main}]\label{thm:main}
For any choice (1), choose \(s>0\) sufficiently small for both retained
tube identifications. The two actual prepared functions, in the analytic
coordinates constructed below, are reduced convergent germs
\[
f_5,f_4:(\C^4,0)\longrightarrow(\C,0)
\]
with isolated critical points, orders five and four, respectively, and
homeomorphic ambient zero-set germs.
\end{theorem}

There is a concrete admissible integer tuple:
\[
r=4,\quad n_0=97,\quad n_1=177,\quad a=2352,\quad b=7832.             \tag{2}
\]
Only the integers are explicit here. No unproved numerical value of the
common analytic smallness threshold for \(s\) is asserted.
For all (1), \(a<b<4a\), \(a>12\), and
\[
\frac{ab-2b}{a+b}>\frac a2-2>4,\qquad
\frac{ab-2a}{a+b}>\frac a2-2>4.
\]
Thus the strict tube conditions hold, and
\(ab/(a+b)>a/2>5\) gives the retained local-isolation condition.
Also \(a,b>4\), the exact condition needed for the initial-form
calculation. The common positive \(s\) exists by taking the smaller of
the two positive thresholds after fixing the integer tuple.

\section{What the cited topological results actually say}

We inspected Kato's original scan \cite{Kato}, not a search summary.
On printed p.~454, Theorem B has \(n\geq3\), not \(n>3\).
In particular \(2n+1=7\) is included. The category is oriented smooth
manifolds and orientation-preserving diffeomorphisms.
The congruence footnote requires a unimodular integral matrix:
\[
A=P^{\mathsf t}BP,\qquad P\in\operatorname{GL}_m(\Z).               \tag{3}
\]
Theorem B concerns two \emph{simple spinnable structures on the standard
sphere} \(S^{2n+1}\). It does not impose a homotopy-sphere hypothesis
on their axes (bindings).

Section 2, printed p.~455, defines a spinnable structure
\(\mathcal S=\{F,h,g\}\): \(F\) is compact with boundary,
\(h\) is a diffeomorphism equal to the identity on an open neighborhood
of \(\partial F\), and \(g:T(F,h)\to M\) is a diffeomorphism.
The quotient \(T(F,h)\) identifies \((x,1)\) with \((h(x),0)\)
and collapses each \(\{y\}\times[0,1]\), \(y\in\partial F\).
The collar coordinate is the positive angular coordinate
\(\exp(2\pi i t)\). On an \(m\)-dimensional ambient manifold, simplicity
means that \(F\) is obtained from a ball by handles of indices at most
\(\lfloor m/2\rfloor\). The inequality in the scan is inclusive;
the OCR text alone is unreliable at this symbol.
An isomorphism carries \(g(F\times t)\) to \(g'(F'\times t)\) for
every \(t\), and hence carries the embedded bindings.

The local analytic example in N\'emethi \cite[pp.~63--64]{Nemethi}
explicitly states that an isolated analytic germ in \(\C^N\) gives
a \emph{simple} spinnable Milnor structure on a small \(S^{2N-1}\).
The definition of simple on the same printed p.~63 is the handle
definition, not a homotopy-bouquet definition. This passage is before
the paper's discussion of polynomial maps at infinity: its later
condition (C) at infinity is irrelevant here. With N\'emethi's
notation \(N=4\), the handle bound is \(N-1=3\); with Kato's notation
the same number is \(n=3\). This avoids a one-dimensional shift in
applying the theorem.

Finally, Massey \cite[Theorem 1.4, p.~4]{Massey} states the analytic
ambient-pair cone theorem for a non-locally-constant analytic function
near a point of its zero set. It includes the ambient ball, its
hypersurface, and the marked center, and is relative to the bounding
sphere. It is not merely an intrinsic conical description of the
hypersurface.

The primary Kato and N\'emethi PDFs and text extractions are retained
in \path{research/m45_topology/literature/}. Massey's versioned
primary text is reused in place at
\path{research/m45_seifert/massey_1410.3312v2.txt}.
Short, explicitly located excerpts accompany this report.
The associated ledger assumptions are
\texttt{A059TOPO-KATO}, \texttt{A059TOPO-MILNOR-SIMPLE}, and
\texttt{A059TOPO-CONE}. They await independent attribution; inspection
by their filer is not passed off as an independent citation check.

\section{Analytic charts and ordinary orders, with the units retained}

Use the manuscript's exact polynomials
\[
\begin{split}
h(X,T)&=T(X-T^2)(X-2T^2),\\
q_1(u,v)&=(u-v-v^8)(u-v-2v^8)(u-v^2),\qquad q_2=q_1-u^2,\\
\psi(X,v)&=(X+v)(X+v-v^2)q_2(X+v,v).
\end{split}                                                       \tag{4}
\]
Write the actual Weierstrass preparations as
\[
G_0=X-T^2-q_1(T,v)=U_0(X,T,v)R_0(X,T,v),\qquad
G_1=T-\psi(X,v)=U_1(X,T,v)R_1(X,T,v).                              \tag{5}
\]
The \(R_i\) are monic in \(v\) of degree four.
The exact specializations, not just leading-term guesses, are
\[
\begin{split}
G_0(0,0,v)&=v^4(1+v^7)(1+2v^7),\\
G_1(0,0,v)&=v^4(1-v)\bigl(1-2v^{15}(1-v)\bigr).
\end{split}                                                       \tag{6}
\]
Consequently \(R_i(0,0,v)=v^4\) and \(U_i(0,0,0)=1\).
Since \(R_i(0)=0\), differentiating (5) at zero gives
\[
\partial_XR_0(0)=1,\qquad \partial_TR_1(0)=1.                      \tag{7}
\]
Put \(R_{i,s}(X,T,w)=s^4R_i(X,T,w/s)\), for fixed \(s\ne0\), and
\[
C_{i,s}=\{yz=R_{i,s}(X,T,w)\},\qquad
f_{i,s}=\left.h(X,T)+w(y^a+z^b)\right|_{C_{i,s}}.
\]
The indicated derivatives of \(yz-R_{i,s}\) at zero are \(-s^4\),
so the holomorphic implicit-function theorem gives smooth ambient
fourfold germs. Solve for \(X\) on \(C_{0,s}\), using
\((T,y,z,w)\) as coordinates, and for \(T\) on \(C_{1,s}\), using
\((X,y,z,w)\).

Set \(c=s^{-4}\ne0\). Multiplying the ambient equation by the
\emph{actual} unit in (5), rather than replacing that unit by one,
gives the exact implicit identities
\[
\begin{split}
X&=T^2+q_1(T,w/s)+c\,U_0(X,T,w/s)yz,\\
T&=\psi(X,w/s)+c\,U_1(X,T,w/s)yz. 
\end{split}                                                       \tag{8}
\]
No function \(f_{i,s}\) or ambient equation has been replaced by an
unprepared polynomial model.

\begin{lemma}[The two initial forms]\label{lem:orders}
In these analytic coordinates,
\[
\begin{split}
X&=T^2+c\,yz+O_{\geq3}\quad\text{on }C_{0,s},\\
T&=c\,yz+O_{\geq3}\quad\text{on }C_{1,s},\\
\operatorname{in}(f_{0,s})&=T(c\,yz)(c\,yz-T^2),\\
\operatorname{in}(f_{1,s})&=c\,yzX^2.
\end{split}                                                       \tag{9}
\]
Thus their ordinary orders are five and four.
\end{lemma}

\begin{proof}
The first nonzero homogeneous part of \(q_1(u,v)\) is
\(u(u-v)^2\), of degree three. The first nonzero homogeneous part of
\(\psi(X,v)\) is \(-(X+v)^4\), of degree four.
The solutions in (8) have no linear part: each right-hand side has
order at least two, even before this observation is substituted back.
Each composed unit is \(1+O_{\geq1}\). Therefore its product with
\(c\,yz\) is \(c\,yz+O_{\geq3}\), proving the two quadratic jets.
The remainder constants may depend on the fixed nonzero \(s\);
no uniform Taylor estimate as \(s\to0\) is needed.

In the first chart both \(X-T^2\) and \(X-2T^2\) have order two, with
initial forms \(c\,yz\) and \(c\,yz-T^2\). Multiplying their initial
forms and \(T\) proves the degree-five formula. In the second chart,
\[
h(X,T)=TX^2-3T^3X+2T^5,
\]
and \(T\) has order two. Only \(TX^2\) contributes in degree four;
the other two displayed terms have orders at least seven and ten.
In both charts \(w(y^a+z^b)\) has order at least six since \(a,b>4\).
It cannot affect either initial form.
For \(c\ne0\), the monomials \(c^2Ty^2z^2\) and \(-cT^3yz\)
are distinct and have nonzero coefficients; \(cyzX^2\) is also
nonzero. Thus these are actual orders, not just lower bounds.
\end{proof}

The retained isolation theorem applies in the range (1).
Reducedness follows from isolation, not from squarefreeness of an
initial form. Indeed, if \(f=g^k u\) for an irreducible nonunit \(g\)
and \(k\geq2\), every first derivative of \(f\) vanishes along \(V(g)\).
That complex hypersurface has positive dimension and nonzero points
arbitrarily close to zero. This contradicts isolation.
In particular the repeated \(X\)-factor in the degree-four tangent
cone is no objection to reducedness of the full analytic germ.

\section{The actual simple open books on standard spheres}

\begin{proposition}\label{prop:simple}
The actual germs in the two charts have simple smooth Milnor open
books on standard oriented \(S^7\)'s, with six-dimensional compact
pages and characteristic diffeomorphisms equal to the identity on
boundary collars.
\end{proposition}

\begin{proof}
Choose an ordinary sufficiently small Euclidean ball in each analytic
chart from (8), not the boundary of the custom tube used to compute
the lattice. Its boundary is literally a round \(S^7\), up to a
positive radial rescaling. A holomorphic chart preserves the complex
orientation: its real Jacobian determinant is the positive number
\(\lvert\det_{\C}D\chi\rvert^2\).
The origin is an isolated critical point by the retained theorem.
Thus N\'emethi's local analytic Milnor statement applies with \(N=4\).
It supplies the compact open-book structure and its handle
simplicity on this standard sphere. Its binding has real dimension
five. Nothing asserts that this binding is a sphere.

We spell out the geometric source of the handle bound, to avoid the
invalid inference from a bouquet homotopy type to handles.
The ordinary compact tube fiber
\[
F_\delta=f^{-1}(\delta)\cap\overline B_\epsilon
\]
is a smooth complex threefold cut out by a regular radius level.
On it, \(\rho(z)=\lVert z\rVert^2\) is strictly plurisubharmonic.
It has no boundary critical points and its gradient points outward
at the boundary. A sufficiently small smooth perturbation, supported
off a boundary collar, is still strictly plurisubharmonic and can
be chosen Morse, with finitely many interior critical points.

At a critical point let \(H\) be its real Hessian and \(J\) the
complex structure. Strict plurisubharmonicity gives
\[
H(v,v)+H(Jv,Jv)>0\quad(v\ne0).
\]
If a negative subspace had real dimension greater than three,
it would meet its \(J\)-image nontrivially in the six-dimensional
tangent space. A nonzero \(v\) and \(Jv\) would then both lie in that
negative subspace, a contradiction. Every Morse index is therefore
at most three. The Morse local normal form and the gradient flow
give actual handles of those indices, not just cells in a
homotopy-equivalent space.
The fiber is connected. If this decomposition begins with several
zero-handles, its one-handles connect them: higher-index handles
cannot join components. A spanning tree of these zero- and
one-handles is a single ball. Absorbing that tree into the initial
ball leaves only handles of indices at most three.

The usual Milnor tube-to-sphere transport, together with the local
product at the binding, transfers this compact page to the compact
open-book page. This is the compact simple structure in the cited
local analytic statement; we are not claiming that an arbitrary
diffeomorphism between interiors extends to arbitrary compactifications.
Thus the preceding handle argument concerns the actual pages.

For completeness, near the binding the open book has the smooth
model \(K\times D^2\), with phase the argument in the disk.
On a smaller such neighborhood choose the angular field to be
the product rotation. On the compact complement, local horizontal
lifts of the unit angular vector field on \(S^1\) can be combined
with a partition of unity, keeping the product field on a still
smaller neighborhood. The lift equation is affine, so the combined
field is still a lift. Its one-turn return fixes a collar of the
page boundary pointwise. It gives the \(h\) required in Kato's
definition, not just a map acting trivially on boundary homology.
\end{proof}

No cancellation of all low-index handles, and no assertion that the
page has only index-three handles, is needed. The bound ``at most
three'' is exactly the one in Kato's definition for ambient dimension
seven.

\section{The ordered form, an onto congruence, and the embedded links}

Let \(L_i=H_3(F_i;\Z)\), with \(F_i\) the pages just constructed.
The retained tube theorem identifies \(L_i\) with its actual lattice
\(\Lambda_i\), and identifies its form with the ordered intersection
of two four-dimensional fillings in the acyclic total tube, the
second boundary on the positive neighboring phase.

\begin{lemma}[Comparison with Kato's convention]\label{lem:convention}
With complex ambient orientation and increasing target argument,
this ordered filling form agrees with Kato's ordered linking form.
In particular the retained integral isometry induces the congruence
(3) in Kato's convention.
\end{lemma}

\begin{proof}
Kato's definition on printed p.~456 is
\[
\Gamma(\alpha,\beta)
 =\lk\bigl(g_*(\alpha\times t_0),g_*(\beta\times t_1)\bigr),
\quad 0\leq t_0<1/2<t_1<1.
\]
The linking number is intersection of a bounding chain for the
first cycle with the second cycle, in that order.
The second page is therefore a \emph{positive} displacement from
the first. A small positive displacement may be increased to
Kato's chosen displacement without crossing the first page or
completing a full turn. This does not change the linking number
or introduce a monodromy factor.

Make the ordinary tube-to-sphere comparison simultaneously over a
closed angular interval containing the two phases. The outgoing
transport cylinders have different fixed phases and are disjoint.
They lie outside a smaller total tube, while the initial filling
chains can be taken inside it. Attaching these cylinders to the
fillings expresses the reviewed pairing as intersection of fillings
in the complex-oriented eight-ball with the transported boundary
cycles on its seven-sphere. The cylinders introduce no additional
intersections. This is why using the \emph{actual} local tube
identification, rather than only the abstract group \(\Lambda_i\),
is important here.

Here is the remaining boundary-orientation check. Give the boundary
of the eight-ball the outward-normal-first orientation. Let
\(C^4\) in its boundary bound the first three-cycle, and let \(b^3\)
be the second three-cycle. Compute the eight-ball intersection by
pushing \(C\) slightly inward and adding its boundary collar.
A filling \(B^4\) of \(b\) has collar orientation
\((n_{\rm out},b)\). Thus the ordered intersection orientation is
\[
(C,n_{\rm out},b)=(-1)^4(n_{\rm out},C,b)
                 =(n_{\rm out},C,b).
\]
It is the same sign as the boundary intersection \(C\cdot b\).
The filling-pair sign is \(+1\) in the present dimension. Since the
ball has vanishing homology in the relevant degrees, changing
fillings does not change these numbers. The same holds in the
acyclic ordinary tube. All transports and both analytic charts
respect the common complex and positive-phase orientations.
This proves the claimed convention comparison.

Let \(I:\Lambda_0\to\Lambda_1\) be the retained \emph{onto} integral
isometry. The groups are free abelian, and \(I^{-1}\) is integral.
In integral bases its matrix \(P\) has determinant \(\pm1\), and
\(\Gamma_0=P^{\mathsf t}\Gamma_1P\). This is an integral congruence,
not a rational congruence or a stable equivalence.
Even if a reader uses the uniformly transposed or negated convention
for Seifert forms, the same \(P\) gives
\((\pm\Gamma_0)^{\mathsf t}
 =P^{\mathsf t}(\pm\Gamma_1)^{\mathsf t}P\).
No independent choice of sign on the two germs is being made.
\end{proof}

We can now check every input to Kato's Theorem B: the ambient
manifolds are the standard oriented \(S^7\)'s; \(n=3\geq3\);
the compact pages have real dimension six and the actual handle
bound at most three; the monodromies fix collars; and the actual
ordered integral Seifert matrices satisfy (3).
The conclusion is an orientation-preserving diffeomorphism
\[
\phi:(S^7_{\epsilon_0},K_0)\longrightarrow
      (S^7_{\epsilon_1},K_1),\qquad
K_i=f_{i,s}^{-1}(0)\cap S^7_{\epsilon_i}.                          \tag{10}
\]
Indeed an isomorphism carries every compact page to its
corresponding page, and hence carries its embedded boundary.
This is stronger than an abstract diffeomorphism \(K_0\cong K_1\).
Theorem A (realization), a spherical-knot classification, and a
homotopy-sphere condition on \(K_i\) are not used.

\section{The ambient germ, not merely the link}

\begin{proof}[Proof of Theorem \ref{thm:main}]
Lemma \ref{lem:orders} and the retained isolation theorem give
reduced isolated analytic germs of the stated orders in two smooth
four-dimensional analytic charts. Their orders are positive, so
they are not locally constant and vanish at the chart origins.
Take the radii in (10) small enough also for the analytic
ambient-pair cone theorem.

Write \(c(S,K)\) for the cone pair with its marked vertex, with
radial parameter \(t=0\) at the vertex and \(t=1\) at its boundary.
Massey's Theorem 1.4 gives pair homeomorphisms
\[
\Psi_i:(\overline B_{\epsilon_i},
        \overline B_{\epsilon_i}\cap V(f_{i,s}),0)
       \longrightarrow c(S^7_{\epsilon_i},K_i),
\]
relative to the outer spheres. No ordinary linear radial
invariance of the analytic zero sets is assumed.
Equation (10) induces
\[
c\phi:[u,t]\longmapsto[\phi(u),t].
\]
It is well defined at the vertex and continuous by the quotient
definition of a cone; \(c(\phi^{-1})\) is its inverse. Therefore
\[
\Psi_1^{-1}\circ c\phi\circ\Psi_0
\]
is a homeomorphism of the two ball pairs fixing the marked centers.
Its restriction to the open ball pairs yields a homeomorphism
germ. Finally conjugate by the two holomorphic charts, and name
the order-five function \(f_5=f_{0,s}\) and the order-four function
\(f_4=f_{1,s}\). This is exactly the assertion
\texttt{M45:thm:main}, with the two roles interchanged if necessary.
\end{proof}

\section{What the exact replay certifies}

\texttt{verify.py} uses rational symbolic arithmetic and the retained
bounded Jacobian factorization. It checks (4)--(9), including the
exact specialized units, the absence of a linear term in the
implicit solutions, the independence of their quadratic jets from
arbitrary positive-degree unit corrections, the nonzero initial
forms, the tuple (2), and the eight-dimensional collar permutation.
The Jacobian payload treats \(c\) as an independent indeterminate;
the proof separately imposes \(c=s^{-4}\ne0\).
Its factorization result is replayed by reconstructing the exact
polynomial. Repeated factors in a tangent cone are not used as a
test for non-reducedness of a germ.
The replay requires Python 3.10 or later and SymPy. From the campaign
root, run \texttt{python research/m45\_topology/verify.py} in the
provided Argus environment. Default mode recomputes the checks,
compares the retained \texttt{exact\_certificate.json}, and emits
a concise JSON result. The explicit \texttt{--write-certificate}
option regenerates that certificate.

The certificate distinguishes these computations from:
(i) the geometric proof of actual simplicity and collar fixing;
(ii) the identification of the filling form with a linking form;
(iii) the cited Kato classification and analytic cone theorem; and
(iv) the retained prior results. It does not claim that a computer
checked any of (i)--(iv). It also does not produce an explicit
Seifert isometry or a formula for the homeomorphism.

The improvement over the strongest retained result is the passage
from an onto integral Seifert isometry to the actual embedded link
equivalence and then the ambient analytic pair germs, with the
different orders established in the correct coordinates.
This is an independent verification argument for M45, not evidence
that these conclusions are original to this audit.
The normal host-invoked Reviewer must assess this proof, replay the
supporting exact calculations, and attribute the three new cited
propositions independently. Existing citation-provenance flags on
retained inputs are not silently erased by this submission.
No conclusion about M23, verified novelty, or campaign-final
publishable acceptance is asserted.

\begin{thebibliography}{9}
\bibitem{Kato}
M. Kato, \emph{A classification of simple spinnable structures on a
1-connected Alexander manifold}, J. Math. Soc. Japan 26 (1974),
454--463, Theorem B and Section 2.
\href{https://doi.org/10.2969/jmsj/02630454}
{doi:10.2969/jmsj/02630454}.

\bibitem{Nemethi}
A. N\'emethi, \emph{On the Seifert form at infinity associated with
polynomial maps}, J. Math. Soc. Japan 51 (1999), 63--70;
definition of simple and the local analytic example, pp.~63--64.
\href{https://doi.org/10.2969/jmsj/05110063}
{doi:10.2969/jmsj/05110063}.

\bibitem{Massey}
D. B. Massey, \emph{Non-isolated Hypersurface Singularities and
L\^e Cycles}, revised November 13, 2015,
\href{https://arxiv.org/abs/1410.3312v2}{arXiv:1410.3312v2},
Theorem 1.4, p.~4.
\end{thebibliography}
\end{document}
